\documentclass[../distribution_theory_notes.tex]{subfiles}
\begin{document}
\section{Aula 22 - 11 de Novembro, 2024}
\subsection{Produto de Convolução -- Terceira Parte}
Vimos que, se
\[
 P(D)=\sum_{|\alpha|\leq m}a_\alpha\partial^\alpha
\]
é um operador diferencial parcial linear com coeficientes constantes e $E\in\mathcal D'(\mathbb R^n)$ é uma solução fundamental, isto é, $P(D)E=\delta$, então, para toda $f\in\mathcal C_c^\infty(\mathbb R^n)$,
\[
 u:=E*f
\]
é uma solução $\mathcal C^\infty$ de $P(D)u=f$.

Para demonstrarmos fatos de existência, unicidade e, mais adiante, ampliarmos o conjunto dos termos não homogêneos, precisamos definir o produto de convolução em situações nas quais ambos os fatores sejam distribuições. O caso básico será aquele em que pelo menos um dos fatores possui suporte compacto.

\begin{def*}
 Para $u\in\mathcal E'(\mathbb R^n)$ e $\varphi\in\mathcal C^\infty(\mathbb R^n)$, ou para $u\in\mathcal D'(\mathbb R^n)$ e $\varphi\in\mathcal C_c^\infty(\mathbb R^n)$, definimos
 \[
  (u*\varphi)(x):=\langle u,\tau_x\check\varphi\rangle
  =\langle u(y),\varphi(x-y)\rangle,
  \qquad x\in\mathbb R^n.
 \]
 Aqui $\tau_x\check\varphi(y)=\varphi(x-y)$.
\end{def*}

\begin{lemma*}
 Para $u\in\mathcal E'(\mathbb R^n)$ e $\varphi\in\mathcal C^\infty(\mathbb R^n)$ valem:
 \begin{enumerate}
  \item $\delta*\varphi=\varphi$;
  \item $u*\varphi\in\mathcal C^\infty(\mathbb R^n)$ e
  \[
   \partial^\alpha(u*\varphi)
   =(\partial^\alpha u)*\varphi
   =u*(\partial^\alpha\varphi),
   \qquad \alpha\in\mathbb Z_+^n;
  \]
  \item
  \[
   (u*\varphi)*\psi=u*(\varphi*\psi),
  \]
  sempre que $\varphi$ ou $\psi$ tiver suporte compacto.
 \end{enumerate}
\end{lemma*}

\begin{proof*}
 As duas primeiras propriedades são as mesmas já demonstradas para o caso $u\in\mathcal D'$ e $\varphi\in\mathcal C_c^\infty$. Para a primeira, escolhendo uma função de Urysohn $\chi\equiv1$ numa vizinhança de $\{0\}$, a extensão de $\delta$ a $\mathcal C^\infty$ satisfaz
 \[
  \langle\delta,f\rangle=\langle\delta,\chi f\rangle=f(0),
 \]
 e, portanto,
 \[
  (\delta*\varphi)(x)=\langle\delta(y),\varphi(x-y)\rangle=\varphi(x).
 \]

 Na associatividade, se, por exemplo, $\psi\in\mathcal C_c^\infty$, então $u*\varphi$ é uma função suave e $(u*\varphi)*\psi$ está bem definido; ao mesmo tempo $\varphi*\psi\in\mathcal C^\infty$ e $u*(\varphi*\psi)$ também está bem definido. O caso em que $\varphi$ tem suporte compacto é análogo. Os detalhes reduzem-se às identidades já vistas e ao teorema de Fubini. \qedsymbol
\end{proof*}

\begin{def*}
 Dados $u_1,u_2\in\mathcal D'(\mathbb R^n)$, com, digamos, $u_2\in\mathcal E'(\mathbb R^n)$, para toda $\varphi\in\mathcal C_c^\infty(\mathbb R^n)$ ficam bem definidas as funções
 \[
  u_1*(u_2*\varphi),\qquad u_2*(u_1*\varphi),
 \]
 pois $u_2*\varphi\in\mathcal C_c^\infty$ e $u_1*\varphi\in\mathcal C^\infty$.
\end{def*}

Nessas condições, consideremos
\[
 T\varphi:=u_1*(u_2*\varphi),
 \qquad
 S\varphi:=u_2*(u_1*\varphi).
\]

\begin{lemma*}
 Os operadores $T,S:\mathcal C_c^\infty(\mathbb R^n)\to\mathcal C^\infty(\mathbb R^n)$ são lineares, sequencialmente contínuos e invariantes por translação:
 \[
  \tau_h(T\varphi)=T(\tau_h\varphi),
  \qquad
  \tau_h(S\varphi)=S(\tau_h\varphi),
  \qquad h\in\mathbb R^n.
 \]
\end{lemma*}

\begin{proof*}
 A invariância por translação decorre das identidades já provadas para a convolução de uma distribuição com uma função teste. Por exemplo,
 \begin{align*}
  \tau_h(T\varphi)
  &=\tau_h\bigl[u_1*(u_2*\varphi)\bigr]\\
  &=u_1*\tau_h(u_2*\varphi)\\
  &=u_1*\bigl[u_2*(\tau_h\varphi)\bigr]
  =T(\tau_h\varphi).
 \end{align*}
 A linearidade é imediata da linearidade das distribuições e das operações de translação e reflexão.

 Para a continuidade, basta mostrar que, se $\varphi_j\to0$ em $\mathcal C_c^\infty$, então $u*\varphi_j\to0$ em $\mathcal C^\infty$ quando $u\in\mathcal E'$. De fato, existe um compacto $K_1$ contendo todos os suportes de $\varphi_j$ e
 \[
  \operatorname{supp}(u*\varphi_j)
  \subseteq\operatorname{supp}(u)+K_1,
 \]
 um compacto fixo. Se
 \[
  |\langle u,f\rangle|
  \leq c\sum_{|\alpha|\leq m}\sup_{y\in K}|\partial^\alpha f(y)|,
 \]
 então, para $x\in\mathbb R^n$,
 \begin{align*}
  |(u*\varphi_j)(x)|
  &=|\langle u,\tau_x\check\varphi_j\rangle|\\
  &\leq c\sum_{|\alpha|\leq m}
  \sup_{z\in\mathbb R^n}|\partial^\alpha\varphi_j(z)|
  \longrightarrow0.
 \end{align*}
 Analogamente, para todo multi-índice $\beta$,
 \[
  |\partial^\beta(u*\varphi_j)(x)|
  =|u*(\partial^\beta\varphi_j)(x)|
  \leq c\sum_{|\alpha|\leq m}
  \sup_z|\partial^{\alpha+\beta}\varphi_j(z)|
  \longrightarrow0.
 \]
 Portanto $u*\varphi_j\to0$ em $\mathcal C^\infty$. Aplicando isso aos dois fatores, obtemos a continuidade de $T$ e $S$. \qedsymbol
\end{proof*}

Pelo resultado que caracteriza operadores lineares contínuos invariantes por translação, existem únicas distribuições $u,v\in\mathcal D'(\mathbb R^n)$ tais que
\[
 u_1*(u_2*\varphi)=u*\varphi,
 \qquad
 u_2*(u_1*\varphi)=v*\varphi,
 \qquad \forall\varphi\in\mathcal C_c^\infty.
\]

\begin{def*}
 Nessas condições definimos
 \[
  u_1*u_2:=u,
  \qquad
  u_2*u_1:=v.
 \]
 Assim,
 \[
  (u_1*u_2)*\varphi=u_1*(u_2*\varphi),
  \qquad
  (u_2*u_1)*\varphi=u_2*(u_1*\varphi).
 \]
\end{def*}

\subsection{Propriedades da Convolução de Distribuições}
Supondo sempre que as convoluções estejam definidas, valem as propriedades usuais.

\begin{prop*}
 \begin{enumerate}
  \item $\delta*u=u*\delta=u$;
  \item $u_1*(u_2*u_3)=(u_1*u_2)*u_3$ se ao menos dois dos fatores tiverem suporte compacto;
  \item $u_1*u_2=u_2*u_1$;
  \item
  \[
   \partial^\alpha(u_1*u_2)
   =(\partial^\alpha u_1)*u_2
   =u_1*(\partial^\alpha u_2),
   \qquad \alpha\in\mathbb Z_+^n;
  \]
  \item
  \[
   \operatorname{supp}(u_1*u_2)
   \subseteq\operatorname{supp}(u_1)+\operatorname{supp}(u_2);
  \]
  \item se um dos fatores é uma função suave, as definições nova e antiga coincidem.
 \end{enumerate}
\end{prop*}

\begin{proof*}
 Para o elemento neutro, se $\varphi\in\mathcal C_c^\infty$,
 \[
  (\delta*u)*\varphi=\delta*(u*\varphi)=u*\varphi,
 \]
 logo $\delta*u=u$; analogamente, $(u*\delta)*\varphi=u*(\delta*\varphi)=u*\varphi$.

 Para a associatividade, convolvemos ambos os lados com uma função teste e usamos repetidamente a própria definição. Por exemplo,
 \begin{align*}
  [u_1*(u_2*u_3)]*\varphi
  &=u_1*[(u_2*u_3)*\varphi]\\
  &=u_1*[u_2*(u_3*\varphi)]\\
  &=(u_1*u_2)*(u_3*\varphi)\\
  &=[(u_1*u_2)*u_3]*\varphi.
 \end{align*}
 Como duas distribuições que têm a mesma convolução com toda função teste coincidem, segue a igualdade.

 Para a comutatividade, tomemos $\varphi,\psi\in\mathcal C_c^\infty$. Pela associatividade já provada,
 \begin{align*}
  (u_1*u_2)*(\varphi*\psi)
  &=(u_1*\psi)*(u_2*\varphi),\\
  (u_2*u_1)*(\varphi*\psi)
  &=(u_2*\varphi)*(u_1*\psi).
 \end{align*}
 Os lados direitos são iguais pela comutatividade usual da convolução de funções. Fazendo $\psi*\varphi_\varepsilon\to\psi$ em $\mathcal C_c^\infty$, concluímos
 \[
  (u_1*u_2)*\psi=(u_2*u_1)*\psi,
 \]
 para toda $\psi$, e portanto $u_1*u_2=u_2*u_1$.

 Para as derivadas, por definição,
 \begin{align*}
  [ (\partial^\alpha u_1)*u_2]*\varphi
  &=(\partial^\alpha u_1)*(u_2*\varphi)\\
  &=u_1*[\partial^\alpha(u_2*\varphi)]\\
  &=[u_1*(\partial^\alpha u_2)]*\varphi,
 \end{align*}
 e a outra igualdade segue tomando uma família regularizante e passando ao limite.

 Para o suporte, usaremos o lema: se $u_j\to u$ em $\mathcal D'$ e $\operatorname{supp}(u_j)\subseteq F_j$, com $F_j$ fechados e $F_j\downarrow F$, então
 \[
  \operatorname{supp}(u)\subseteq F.
 \]
 De fato, se $\varphi\in\mathcal C_c^\infty(\mathbb R^n\setminus F)$, para $j$ grande seu suporte é disjunto de $F_j$, logo $\langle u_j,\varphi\rangle=0$ e, passando ao limite, $\langle u,\varphi\rangle=0$.

 Agora, se $\varphi_\varepsilon$ é uma família regularizante,
 \[
  (u_1*u_2)*\varphi_\varepsilon
  =u_1*(u_2*\varphi_\varepsilon),
 \]
 e, pelo lema do suporte para convolução de distribuição com função,
 \begin{align*}
  \operatorname{supp}[(u_1*u_2)*\varphi_\varepsilon]
  &\subseteq\operatorname{supp}(u_1)
  +\operatorname{supp}(u_2*\varphi_\varepsilon)\\
  &\subseteq\operatorname{supp}(u_1)+\operatorname{supp}(u_2)
  +\overline{B(0,\varepsilon)}.
 \end{align*}
 Como $(u_1*u_2)*\varphi_\varepsilon\to u_1*u_2$ em $\mathcal D'$, o lema anterior dá
 \[
  \operatorname{supp}(u_1*u_2)
  \subseteq\operatorname{supp}(u_1)+\operatorname{supp}(u_2).
 \]

 Finalmente, quando $u_1$ ou $u_2$ é uma função suave, a associatividade já demonstrada para uma distribuição e funções mostra que a definição por operadores coincide com a definição anterior. \qedsymbol
\end{proof*}

\end{document}
